\ifdefined\gkver\else\def\gkver{student}\fi
\documentclass[\gkver]{gaokaozhenti}
\begin{document}

\chapter{2019年全国理综Ⅱ}
%% paperType: 理综物理部分

\begin{enumerate}
\item
%% number: 14
%% paperName: 2019年全国理综Ⅱ第 14 题 6 分
%% typeId: 1
%% type: 单选题
%% chapter: 曲线运动
%% point: 万有引力定律
%% method: 无
%% score: 6
%% degree: 750
%% duplicateId: 0
%% body:
2019年1月，我国嫦娥四号探测器成功在月球背面软着陆，在探测器“奔向”月球的过程中，用 $h$ 表示探测器与地球表面的距离，$F$ 表示它所受的地球引力，能够描述 $F$ 随 $h$ 变化关系的图像是\xzanswer{D}

\fourchoices[answer=D, ispicture=true, h=2.4cm]
{figs/14a.png}
{figs/14b.png}
{figs/14c.png}
{figs/14d.png}

%% answer: D
%% memo:
\memoanswer{
设地球半径为 $R$，地球质量为 $M$，探测器质量为 $m$，探测器在远离地球的过程中，其所受地球的引力大小为 $F=\frac{GMm}{(R+h)^2}$，由此可知，随着 $h$ 的增大，$F$ 逐渐减小，但 $F$ 随 $h$ 不是均匀减小的，故 D 正确。
}

\item
%% number: 15
%% paperName: 2019年全国理综Ⅱ第 15 题 6 分
%% typeId: 1
%% type: 单选题
%% chapter: 原子和原子核
%% point: 质能方程
%% method: 无
%% score: 6
%% degree: 750
%% duplicateId: 0
%% body:
太阳内部核反应的主要模式之一是质子-质子循环，循环的结果可表示为 $4\ce{^1_1H}\to\ce{^4_2He} + 2\ce{^0_1e} + 2\nu$，已知 $\ce{^1_1H}$ 和 $\ce{^4_2He}$ 的质量分别为 $m_{\mathrm{P}}=1.0078\Uu$ 和 $m_{\alpha}=4.0026\Uu$，$1\Uu=931\UMeVcSq$，$c$ 为光速。在 4 个 $\ce{^1_1H}$ 转变成 1 个 $\ce{^4_2He}$ 的过程中，释放的能量约为\xzanswer{C}
\fourchoices[answer=C]
{$8\UMeV$}
{$16\UMeV$}
{$26\UMeV$}
{$52\UMeV$}
%% answer: C
%% memo:
\memoanswer{
该核反应过程的质量亏损约为 $\Delta m=4m_{\mathrm{p}}-m_{\alpha}=0.0286\Uu$，得该反应释放的能量约为 $\Delta E=0.0286\times931\UMeV=26.63\UMeV$，故 C 正确。
}

\item
%% number: 16
%% paperName: 2019年全国理综Ⅱ第 16 题 6 分
%% typeId: 1
%% type: 单选题
%% chapter: 力物体的平衡
%% point: 多力平衡
%% method: 无
%% score: 6
%% degree: 650
%% duplicateId: 0
%% body:
物块在轻绳的拉动下沿倾角为 $30^{\circ}$ 的固定斜面向上匀速运动，轻绳与斜面平行。已知物块与斜面之间的动摩擦因数为 $\frac{\sqrt{3}}{3}$，重力加速度取 $10\Umsq$。若轻绳能承受的最大张力为 $1500\UN$，则物块的质量最大为\xzanswer{A}
\fourchoices[answer=A]
{$150\Ukg$}
{$100\sqrt{3}\Ukg$}
{$200\Ukg$}
{$200\sqrt{3}\Ukg$}
%% answer: A
%% memo:
\memoanswer{
对物块受力分析如图所示，物块受重力 $mg$，沿斜面向下的摩擦力 $f$，垂直于斜面向上的支持力 $N$ 和绳的拉力 $T$。设物块的最大质量为 $m$，垂直于斜面方向有 $N=mg\cos30^{\circ}$，沿斜面方向有 $T=mg\sin30^{\circ}+f$，$f=\mu N$，联立以上三式解得 $m=150\Ukg$，故 A 正确。

\onepicture[w=4.5cm]{figs/16_an.png}
}

\item
%% number: 17
%% paperName: 2019年全国理综Ⅱ第 17 题 6 分
%% typeId: 1
%% type: 单选题
%% chapter: 磁场
%% point: 带电粒子在磁场中的运动
%% method: 无
%% score: 6
%% degree: 550
%% duplicateId: 0
%% body:
如图，边长为 $l$ 的正方形 abcd 内存在匀强磁场，磁感应强度大小为 $B$，方向垂直于纸面（abcd 所在平面）向外。ab 边中点有一电子发源 O，可向磁场内沿垂直于 ab 边的方向发射电子。已知电子的比荷为 $k$。则从 a、d 两点射出的电子的速度大小分别为\xzanswer{B}
\onepicture[w=5.5cm, align=right]{figs/17.png}
\fourchoices[answer=B]
{$\frac{1}{4}kBl$，$\frac{\sqrt{5}}{4}kBl$}
{$\frac{1}{4}kBl$，$\frac{5}{4}kBl$}
{$\frac{1}{2}kBl$，$\frac{\sqrt{5}}{4}kBl$}
{$\frac{1}{2}kBl$，$\frac{5}{4}kBl$}
%% answer: B
%% memo:
\memoanswer{
电子在匀强磁场中运动，洛伦兹力提供其做圆周运动的向心力，由 $qvB=\frac{mv^2}{r}$ 得，电子在匀强磁场中运动的轨迹半径为 $r=\frac{mv}{qB}$。若电子从 a 点射出，则轨迹圆心在 Oa 中点上，Oa 之间的轨迹为半圆，则轨迹半径 $r_a=\frac{l}{4}$，解得从 a 点射出的电子速度大小 $v_a=\frac{qBl}{4m}=\frac{kBl}{4}$。若电子从 d 点射出，运动轨迹如图所示，设轨迹半径为 $r_d$，则 $r_d^2=l^2+\left(r_d-\frac{l}{2}\right)^2$，解得 $r_d=\frac{5}{4}l$，与 $r=\frac{mv}{qB}$ 联立解得从 d 点射出的电子速度大小为 $v_d=\frac{5qBl}{4m}=\frac{5kBl}{4}$，故 B 正确。

\onepicture[w=5cm]{figs/17_an.png}
}

\item
%% number: 18
%% paperName: 2019年全国理综Ⅱ第 18 题 6 分
%% typeId: 2
%% type: 多选题
%% chapter: 功和能
%% point: 功能中的图象
%% method: 无
%% score: 6
%% degree: 450
%% duplicateId: 0
%% body:
从地面竖直向上抛出一物体，其机械能 $E_{\text{总}}$ 等于动能 $E_k$ 与重力势能 $E_p$ 之和。取地面为重力势能零点，该物体的 $E_{\text{总}}$ 和 $E_p$ 随它离开地面的高度 $h$ 的变化如图所示。重力加速度取 $10\Umsq$。由图中数据可得\xzanswer{AD}
\onepicture[w=5.5cm, align=right]{\includesvg{figs/18}}
\fourchoices[answer=AD]
{物体的质量为 $2\Ukg$}
{$h=0$ 时，物体的速率为 $20\Ums$}
{$h=2\Um$ 时，物体的动能 $E_k=40\UJ$}
{从地面至 $h=4\Um$，物体的动能减少 $100\UJ$}
%% answer: AD
%% memo:
\memoanswer{
由 $\frac{\Delta E_p}{\Delta h}=mg$ 可知 $E_p-h$ 图像的斜率大小表示重力大小，所以物体的质量为 $m=2\Ukg$，故 A 正确；当 $h=0$ 时，物体的重力势能为零，此时机械能大小等于物体动能的大小，由图可知，$\frac{1}{2}mv^2=100\UJ$，代入质量可得物体的速率为 $v=10\Ums$，故 B 错误；由图可知，在 $h=2\Um$ 时，物体的重力势能大小为 $40\UJ$，机械能大小为 $90\UJ$，则物体的动能大小为 $50\UJ$，故 C 错误；在地面时，物体的动能为 $100\UJ$，在 $h=4\Um$ 处，机械能等于重力势能，此时动能为零，所以从地面至 $h=4\Um$ 处，物体的动能减少了 $100\UJ$，故 D 正确。
}

\item
%% number: 19
%% paperName: 2019年全国理综Ⅱ第 19 题 6 分
%% typeId: 2
%% type: 多选题
%% chapter: 曲线运动
%% point: 平抛运动
%% method: 无
%% score: 6
%% degree: 450
%% duplicateId: 0
%% body:
如图（a），在跳台滑雪比赛中，运动员在空中滑翔时身体的姿态会影响其下落的速度和滑翔的距离。某运动员先后两次从同一跳台起跳，每次都从离开跳台开始计时，用 $v$ 表示他在竖直方向的速度，其 $v-t$ 图像如图（b）所示，$t_1$ 和 $t_2$ 是他落在倾斜雪道上的时刻。则\xzanswer{BD}
\onepicture[w=8cm, align=right]{figs/19.png}
\fourchoices[answer=BD]
{第二次滑翔过程中在竖直方向上的位移比第一次的小}
{第二次滑翔过程中在水平方向上的位移比第一次的大}
{第二次滑翔过程中在竖直方向上的平均加速度比第一次的大}
{竖直方向速度大小为 $v_1$ 时，第二次滑翔在竖直方向上所受阻力比第一次的大}
%% answer: BD
%% memo:
\memoanswer{
由题图可知该图线的斜率表示运动员在竖直方向的加速度大小，所以两次平均加速度大小满足 $a_1>a_2$，故 C 错误；由 $mg-f=ma$，两次竖直方向速度为 $v_1$ 时有第一次滑翔在竖直方向的加速度大小大于第二次的，可知第一次在竖直方向受到的阻力小于第二次的，故 D 正确；$v-t$ 图像中图线与时间轴围成的图形的面积表示位移的大小，所以由题图可知第二次滑翔过程的竖直分位移大小比第一次的大，故 A 错误；设倾斜雪道的倾角为 $\theta$，则 $\tan\theta=\frac{h}{x}$，所以 $x=\frac{h}{\tan\theta}$，$h_2>h_1$，所以第二次滑翔在水平方向的位移比第一次的大，故 B 正确。
}

\item
%% number: 20
%% paperName: 2019年全国理综Ⅱ第 20 题 6 分
%% typeId: 2
%% type: 多选题
%% chapter: 电场
%% point: 带电粒子的往复运动
%% method: 无
%% score: 6
%% degree: 450
%% duplicateId: 0
%% body:
静电场中，一带电粒子仅在电场力的作用下自 M 点由静止开始运动，N 为粒子运动轨迹上的另外一点，则\xzanswer{AC}
\fourchoices[answer=AC]
{运动过程中，粒子的速度大小可能先增大后减小}
{在 M、N 两点间，粒子的轨迹一定与某条电场线重合}
{粒子在 M 点的电势能不低于其在 N 点的电势能}
{粒子在 N 点所受电场力的方向一定与粒子轨迹在该点的切线平行}
%% answer: AC
%% memo:
\memoanswer{
带电粒子仅在电场力作用下由 M 点运动到 N 点，由于电场线不一定是直线，运动轨迹可能是直线，也可能是曲线，所以粒子的速度可能一直增大，也可能先增大后减小，故 A 正确；若电场线是曲线，则带电粒子的受力不能与其速度方向时刻保持一致，粒子的运动轨迹不会和电场线重合，故 B 错误；只有电场力做功，由功能关系可知，粒子动能和电势能的总和不变，粒子在 M 点的动能为零，可知粒子在 M 点的电势能最大，故 C 正确；运动轨迹的切线方向是速度的方向，而粒子受到的电场力方向与该点电场线的切线方向平行，故 D 错误。
}

\item
%% number: 21
%% paperName: 2019年全国理综Ⅱ第 21 题 6 分
%% typeId: 2
%% type: 多选题
%% chapter: 电磁感应
%% point: 电磁感应中图象
%% method: 无
%% score: 6
%% degree: 350
%% duplicateId: 0
%% body:
如图，两条光滑平行金属导轨固定，所在平面与水平面夹角为 $\theta$，导轨电阻忽略不计。虚线 ab、cd 均与导轨垂直，在 ab 与 cd 之间的区域存在垂直于导轨所在平面的匀强磁场。将两根相同的导体棒 PQ、MN 先后自导轨上同一位置由静止释放，两者始终与导轨垂直且接触良好。已知 PQ 进入磁场开始计时，到 MN 离开磁场区域为止，流过 PQ 的电流随时间变化的图像可能正确的是\xzanswer{AD}
\onepicture[w=7cm, align=right]{\includesvg{figs/21}}

\fourchoices[answer=AD, ispicture=true, h=2.2cm]
{figs/21a.png}
{figs/21b.png}
{figs/21c.png}
{figs/21d.png}

%% answer: AD
%% memo:
\memoanswer{
两导体棒从同一位置由静止释放，进入磁场的速度大小相同，设为 $v$，PQ 刚进入磁场时刚好做匀速运动，则 $mg\sin\theta=F$，$F=BI_1l$，$I_1=\frac{Blv}{R}$，则 $mg\sin\theta=\frac{B^2l^2v}{R}$。从 PQ 进磁场到 MN 出磁场的过程中，可能的情况有两种：第一种情况是 PQ 在磁场中运动时，MN 还没有进入磁场，等 PQ 出了磁场后，MN 再进入磁场，这种情况 PQ 和 MN 分别都是在磁场中做匀速直线运动，产生的感应电流大小相等，但是回路中电流的方向相反，故 A 正确 B 错误；第二种情况是 PQ 出磁场之前，MN 已经进入磁场，MN 刚进磁场时，由于两根导体棒速度大小相同，回路磁通量变化率为零，即感应电动势为零，电流为零，两棒均以 $a=g\sin\theta$ 的加速度做加速运动，当 PQ 出磁场后，MN 在磁场中切割磁感线，产生感应电流，此时切割磁感线的速度大于刚进入磁场时的速度，产生的电流大于 PQ 刚进磁场时的电流，MN 受到的合外力沿导轨向上，$BI'l-mg\sin\theta=ma$，$I'=\frac{Blv'}{R}$，MN 做加速度减小的减速运动，PQ 中的感应电流大小逐渐减小，方向与之前相反，$I-t$ 图像的斜率逐渐减小，故 C 错误 D 正确。
}

\item
%% number: 22
%% paperName: 2019年全国理综Ⅱ第 22 题 5 分
%% typeId: 4
%% type: 实验题
%% chapter: 直线运动
%% point: 打点计时器公式
%% method: 无
%% score: 5
%% degree: 550
%% duplicateId: 0
%% body:
如图（a），某同学设计了测量铁块与木板间动摩擦因数的实验。所用器材有：铁架台、长木板、铁块、米尺、电磁打点计时器、频率 $50\UHz$ 的交流电源，纸带等。回答下列问题：
\begin{enumerate}
	\item
	铁块与木板间动摩擦因数 $\mu=$ \tkanswer[2.2cm]{$\frac{g\sin\theta-a}{g\cos\theta}$}（用木板与水平面的夹角 $\theta$、重力加速度 $g$ 和铁块下滑的加速度 $a$ 表示）
	\item
	某次实验时，调整木板与水平面的夹角 $\theta=30^{\circ}$。接通电源。开启打点计时器，释放铁块，铁块从静止开始沿木板滑下。多次重复后选择点迹清晰的一条纸带，如图（b）所示。图中的点为计数点（每两个相邻的计数点间还有 4 个点未画出）。重力加速度为 $9.8\Umsq$。可以计算出铁块与木板间的动摩擦因数为 \tkanswer[2.2cm]{0.35}（结果保留 2 位小数）。
\end{enumerate}
\twopicture[w=6.5cm, align=right]{figs/22a.png}{figs/22b.png}
%% answer:
\jdanswer{
\begin{enumerate}
	\item
	$\frac{g\sin\theta-a}{g\cos\theta}$

	\item
	0.35

\end{enumerate}
}
%% memo:
\memoanswer{
（1）对铁块受力分析，铁块受重力、支持力和摩擦力，铁块沿木板向下运动，根据牛顿第二定律有 $mg\sin\theta-\mu mg\cos\theta=ma$，解得 $\mu=\frac{g\sin\theta-a}{g\cos\theta}$。

（2）由纸带中的数据，舍弃中央段数据，根据逐差法可求得铁块的加速度大小为 $a=\frac{[(76.39-31.83)-20.90]\times10^{-2}}{3\times4\times0.01}\Umsq=1.97\Umsq$，解得 $\mu=\frac{g\sin\theta-a}{g\cos\theta}=0.35$。
}

\item
%% number: 23
%% paperName: 2019年全国理综Ⅱ第 23 题 10 分
%% typeId: 4
%% type: 实验题
%% chapter: 电路
%% point: 传感器
%% method: 无
%% score: 10
%% degree: 450
%% duplicateId: 0
%% body:
某小组利用图（a）所示的电路，研究硅二极管在恒定电流条件下的正向电压 $U$ 与温度 $t$ 的关系，图中 $V_1$ 和 $V_2$ 为理想电压表；$R$ 为滑动变阻器，$R_0$ 为定值电阻（阻值 $100\UO$）；S 为开关，$E$ 为电源。实验中二极管置于控温炉内，控温炉内的温度 $t$ 由温度计（图中未画出）测出。图（b）是该小组在恒定电流为 $50.0\UuA$ 时得到的某硅二极管 $U-I$ 关系曲线。回答下列问题：
\begin{enumerate}
	\item
	实验中，为保证流过二极管的电流为 $50.0\UuA$，应调节滑动变阻器 $R$，使电压表 $V_1$ 的示数为 $U_1=$ \tkanswer[1.8cm]{5.00} $\UmV$；根据图（b）可知，当控温炉内的温度 $t$ 升高时，硅二极管正向电阻 \tkanswer[1.5cm]{变小}（填“变大”或“变小”），电压表 $V_1$ 示数 \tkanswer[1.5cm]{增大}（填“增大”或“减小”），此时应将 $R$ 的滑片向 \tkanswer[1cm]{B}（填“A”或“B”）端移动，以使 $V_1$ 示数仍为 $U_1$。
	\item
	由图（b）可以看出 $U$ 与 $t$ 成线性关系，硅二极管可以作为测温传感器，该硅二极管的测温灵敏度为 $\left|\frac{\Delta U}{\Delta t}\right|=$ \tkanswer[1.5cm]{2.8} $\times10^{-3}\UV/\UdC$（保留 2 位有效数字）。
\end{enumerate}
\twopicture[w=6.5cm, align=right]{figs/23a.png}{figs/23b.png}
%% answer:
\jdanswer{
\begin{enumerate}
	\item
	5.00，变小，增大，B

	\item
	2.8

\end{enumerate}
}
%% memo:
\memoanswer{
（1）由于题中的两个电压表为理想电表，由图可知，电阻 $R_0$ 与二极管、滑动变阻器串联接入电路，即流过二极管的电流与流过 $R_0$ 的电流相等，又知该电路为恒流电路，可得 $R_0$ 两端的电压 $U_1=IR_0=50.0\UuA\times100\UO=5.00\UmV$，由图（b）可知，当温度升高时，二极管两端的电压减小，而电流不变，由欧姆定律 $R=\frac{U}{I}$ 可知，该二极管的正向电阻变小，由于二极管的正向电阻减小，使得外电路的总电阻减小，电路中的总电流增大，所以电压表 $V_1$ 的示数要增大。由于要保证整个电路中电流不变，即电路中的总电阻不变，可知当温度升高时，需要调整滑动变阻器的滑片向 B 端移动。

（2）该硅二极管的测温灵敏度 $\left|\frac{\Delta U}{\Delta t}\right|=\frac{0.44-0.30}{80-30}\UV/\UdC=2.8\times10^{-3}\UV/\UdC$。
}

\item
%% number: 24
%% paperName: 2019年全国理综Ⅱ第 24 题 12 分
%% typeId: 6
%% type: 计算题
%% chapter: 电场
%% point: 带电粒子的偏转
%% method: 无
%% score: 12
%% degree: 450
%% duplicateId: 0
%% body:
如图，两金属板 P、Q 水平放置，间距为 $d$。两金属板正中间有一水平放置的金属网 G，PQG 的尺寸相同。G 接地，PQ 的电势均为 $\varphi$（$\varphi>0$）。质量为 $m$，电荷量为 $q$（$q>0$）的粒子自 G 的左端上方距离 G 为 $h$ 的位置，以速度 $v_0$ 平行于纸面水平射入电场，重力忽略不计。
\begin{enumerate}
	\item
	求粒子第一次穿过 G 时的动能，以及它从射入电场至此时在水平方向上的位移大小；
	\item
	若粒子恰好从 G 的下方距离 G 也为 $h$ 的位置离开电场，则金属板的长度最短应为多少？
\end{enumerate}
\onepicture[w=6.5cm, align=right]{figs/24.png}
%% answer:
\jdanswer{
\begin{enumerate}
	\item
	$E_k=\frac{1}{2}mv_0^2+\frac{2\varphi}{d}qh$，$l=v_0\sqrt{\frac{mdh}{q\varphi}}$

	\item
	$L=2v_0\sqrt{\frac{mdh}{q\varphi}}$

\end{enumerate}
}
%% memo:
\memoanswer{
（1）PG、QG 间场强大小相等，均为 $E$。粒子在 PG 间所受电场力 $F$ 的方向竖直向下，设粒子的加速度大小为 $a$，则有 $E=\frac{2\varphi}{d}$①，$F=qE=ma$②。设粒子第一次到达 G 时动能为 $E_k$，由动能定理得 $qEh=E_k-\frac{1}{2}mv_0^2$③。设粒子第一次到达 G 时所用的时间为 $t$，粒子在水平方向的位移大小为 $l$，则有 $h=\frac{1}{2}at^2$④，$l=v_0t$⑤。联立①②③④⑤式解得 $E_k=\frac{1}{2}mv_0^2+\frac{2\varphi}{d}qh$⑥，$l=v_0\sqrt{\frac{mdh}{q\varphi}}$⑦。

（2）若粒子穿过 G 一次就从电场的右侧飞出，则金属板的长度最短。由对称性知，此时金属板的长度为 $L=2l=2v_0\sqrt{\frac{mdh}{q\varphi}}$⑧。
}

\item
%% number: 25
%% paperName: 2019年全国理综Ⅱ第 25 题 20 分
%% typeId: 6
%% type: 计算题
%% chapter: 直线运动
%% point: 变加速直线运动
%% method: 无
%% score: 20
%% degree: 250
%% duplicateId: 0
%% body:
一质量为 $m=2000\Ukg$ 的汽车以某一速度在平直公路上匀速行驶。行驶过程中，司机忽然发现前方 $100\Um$ 处有一警示牌。立即刹车。刹车过程中，汽车所受阻力大小随时间变化可简化为图（a）中的图线。图（a）中，$0\sim t_1$ 时间段为从司机发现警示牌到采取措施的反应时间（这段时间内汽车所受阻力已忽略，汽车仍保持匀速行驶），$t_1=0.8\Us$；$t_1\sim t_2$ 时间段为刹车系统的启动时间，$t_2=1.3\Us$；从 $t_2$ 时刻开始汽车的刹车系统稳定工作，直至汽车停止，已知从 $t_2$ 时刻开始，汽车第 $1\Us$ 内的位移为 $24\Um$，第 $4\Us$ 内的位移为 $1\Um$。
\begin{enumerate}
	\item
	在图（b）中定性画出从司机发现警示牌到刹车系统稳定工作后汽车运动的 $v-t$ 图线；
	\item
	求 $t_2$ 时刻汽车的速度大小及此后的加速度大小；
	\item
	求刹车前汽车匀速行驶时的速度大小及 $t_1\sim t_2$ 时间内汽车克服阻力做的功；司机发现警示牌到汽车停止，汽车行驶的距离约为多少（以 $t_1\sim t_2$ 时间段始末速度的算术平均值替代这段时间内汽车的平均速度）？
\end{enumerate}
\onepicture[w=9cm, align=right]{figs/25.png}
%% answer:
\jdanswer{
\begin{enumerate}
	\item
	见解析

	\item
	$a=8\Umsq$，$v_2=28\Ums$

	\item
	$v_1=30\Ums$；$1.16\times10^{5}\UJ$；$87.5\Um$

\end{enumerate}
}
%% memo:
\memoanswer{
（1）$v-t$ 图像如图所示。

\onepicture[w=6cm]{figs/25_an.png}

（2）设刹车前汽车匀速行驶时的速度大小为 $v_1$，则 $t_1$ 时刻的速度也为 $v_1$；$t_2$ 时刻的速度为 $v_2$。在 $t_2$ 时刻后汽车做匀减速运动，设其加速度大小为 $a$，取 $\Delta t=1\Us$，设汽车在 $t_2+(n-1)\Delta t\sim t_2+n\Delta t$ 内的位移为 $s_n(n=1,2,3,\ldots)$。若汽车在 $t_2+3\Delta t\sim t_2+4\Delta t$ 时间内未停止，设它在 $t_2+3\Delta t$ 时刻的速度为 $v_3$，在 $t_2+4\Delta t$ 时刻的速度为 $v_4$，由运动学公式有 $s_1-s_4=3a(\Delta t)^2$①，$s_1=v_2\Delta t-\frac{1}{2}a(\Delta t)^2$②，$v_4=v_2-4a\Delta t$③。联立①②③式，代入已知数据解得 $v_4=-\frac{17}{6}\Ums$④。这说明在 $t_2+4\Delta t$ 时刻前，汽车已经停止。因此，①式不成立。

由于在 $t_2+3\Delta t\sim t_2+4\Delta t$ 内汽车停止，由运动学公式得 $v_3=v_2-3a\Delta t$⑤，$2as_4=v_3^2$⑥。联立②⑤⑥式，代入已知数据解得 $a=8\Umsq$，$v_2=28\Ums$⑦，或 $a=\frac{288}{25}\Umsq$，$v_2=29.76\Ums$⑧。但⑧式情形下，$v_3<0$，不合题意，舍去。

（3）设汽车的刹车系统稳定工作时，汽车所受阻力的大小为 $f_1$，由牛顿第二定律得 $f_1=ma$⑨。在 $t_1\sim t_2$ 时间内，阻力对汽车冲量的大小为 $I=\frac{1}{2}f_1(t_2-t_1)$⑩。由动量定理得 $I=mv_1-mv_2$\textcircled{\scriptsize 11}。由动能定理，在 $t_1\sim t_2$ 时间内，汽车克服阻力做的功为 $W=\frac{1}{2}mv_1^2-\frac{1}{2}mv_2^2$\textcircled{\scriptsize 12}。联立⑦⑨⑩\textcircled{\scriptsize 11}\textcircled{\scriptsize 12}式，代入已知数据解得 $v_1=30\Ums$\textcircled{\scriptsize 13}，$W=1.16\times10^{5}\UJ$\textcircled{\scriptsize 14}。从司机发现警示牌到汽车停止，汽车行驶的距离 $s$ 约为 $s=v_1t_1+\frac{1}{2}(v_1+v_2)(t_2-t_1)+\frac{v_2^2}{2a}$\textcircled{\scriptsize 15}。联立⑦\textcircled{\scriptsize 13}\textcircled{\scriptsize 15}式，代入已知数据解得 $s=87.5\Um$\textcircled{\scriptsize 16}。
}

\item
%% number: 33
%% paperName: 2019年全国理综Ⅱ第 33 题 10 分
%% typeId: 6
%% type: 计算题
%% chapter: 气体性质
%% point: 气体连接体
%% method: 无
%% score: 10
%% degree: 450
%% duplicateId: 0
%% body:
\begin{enumerate}
	\item
	（5 分）如 $p-V$ 图所示，1、2、3 三个点代表某容器中一定量理想气体的三个不同状态，对应的温度分别是 $T_1$、$T_2$、$T_3$。用 $N_1$、$N_2$、$N_3$ 分别表示这三个状态下气体分子在单位时间内撞击容器壁上单位面积的次数，则 $N_1$ \tkanswer[1.2cm]{大于} $N_2$，$T_1$ \tkanswer[1.2cm]{等于} $T_3$，$N_2$ \tkanswer[1.2cm]{大于} $N_3$。（填“大于”“小于”或“等于”）
	\onepicture[w=5cm, align=right]{figs/33a.png}
	\item
	（10 分）如图，一容器由横截面积分别为 $2S$ 和 $S$ 的两个汽缸连通而成，容器平放在地面上，汽缸内壁光滑。整个容器被通过刚性杆连接的两活塞分隔成三部分，分别充有氢气、空气和氮气。平衡时，氮气的压强和体积分别为 $p_0$ 和 $V_0$，氢气的体积为 $2V_0$，空气的压强为 $p$。现缓慢地将中部的空气全部抽出，抽气过程中氢气和氮气的温度保持不变，活塞没有到达两汽缸的连接处，求：
	\begin{enumerate}
		\item
		抽气前氢气的压强；
		\item
		抽气后氢气的压强和体积。
	\end{enumerate}
	\onepicture[w=7cm, align=right]{figs/33b.png}
\end{enumerate}
%% answer:
\jdanswer{
\begin{enumerate}
	\item
	大于，等于，大于

	\item
	$\frac{1}{2}(p_0+p)$，$\frac{1}{2}p_0+\frac{1}{4}p$，$\frac{4(p_0+p)V_0}{2p_0+p}$

\end{enumerate}
}
%% memo:
\memoanswer{
（1）根据理想气体状态方程，对于同一系统有 $\frac{p_1V_1}{T_1}=\frac{p_2V_2}{T_2}=\frac{p_3V_3}{T_3}$，所以 $T_1:T_2:T_3=p_1V_1:p_2V_2:p_3V_3=2:1:2$，所以 $T_1=T_3>T_2$。状态 1 和状态 2 气体的体积相等，由于 $T_1>T_2$，所以 $N_1>N_2$；状态 2 和状态 3 气体的压强相同，由于 $T_3>T_2$，所以 $N_2>N_3$。

（2）（i）设抽气前氢气的压强为 $p_{10}$，由力的平衡条件得 $(p_{10}-p)\cdot2S=(p_0-p)\cdot S$①，解得 $p_{10}=\frac{1}{2}(p_0+p)$②。

（ii）设抽气后氢气的压强和体积分别为 $p_1$ 和 $V_1$，氮气的压强和体积分别为 $p_2$ 和 $V_2$。由力的平衡条件得 $p_2\cdot S=p_1\cdot2S$③，由玻意耳定律得 $p_1V_1=p_{10}\cdot2V_0$④，$p_2V_2=p_0V_0$⑤。由于两活塞用刚性杆连接，故 $V_1-2V_0=2(V_0-V_2)$⑥。联立②③④⑤⑥式解得 $p_1=\frac{1}{2}p_0+\frac{1}{4}p$⑦，$V_1=\frac{4(p_0+p)V_0}{2p_0+p}$⑧。
}

\item
%% number: 34
%% paperName: 2019年全国理综Ⅱ第 34 题 10 分
%% typeId: 4
%% type: 实验题
%% chapter: 光学
%% point: 测定玻璃折射率
%% method: 无
%% score: 10
%% degree: 650
%% duplicateId: 0
%% body:
\begin{enumerate}
	\item
	（5 分）如图，长为 $l$ 的细绳下方悬挂一小球 $a$。绳的另一端固定在天花板上 O 点处，在 O 点正下方 $\frac{3}{4}l$ 的 O$'$ 处有一固定细铁钉。将小球向右拉开，使细绳与竖直方向成一小角度（约为 $2^{\circ}$）后由静止释放，并从释放时开始计时。当小球 $a$ 摆至最低位置时，细绳会受到铁钉的阻挡。设小球相对于其平衡位置的水平位移为 $x$，向右为正。下列图像中，能描述小球在开始一个周期内的 $x-t$ 关系的是\xzanswer{A}
	\onepicture[h=4.5cm, align=right]{figs/34.png}
	\fourchoices[answer=A, ispicture=true, h=2.0cm]
	{figs/34a.png}
	{figs/34b.png}
	{figs/34c.png}
	{figs/34d.png}
	\item
	（10 分）某同学利用图示装置测量某种单色光的波长。实验时，接通电源使光源正常发光：调整光路，使得从目镜中可以观察到干涉条纹。回答下列问题：
	\begin{enumerate}
		\item
		若想增加从目镜中观察到的条纹个数，该同学可 \tkanswer[1cm]{B}；
		\fourchoices[answer=B]
		{将单缝向双缝靠近}
		{将屏向靠近双缝的方向移动}
		{将屏向远离双缝的方向移动}
		{使用间距更小的双缝}
		\item
		若双缝的间距为 $d$，屏与双缝间的距离为 $l$，测得第 1 条暗条纹到第 $n$ 条暗条纹之间的距离为 $\Delta x$，则单色光的波长 $\lambda=$ \tkanswer[2cm]{$\frac{\Delta x\cdot d}{(n-1)l}$}；
		\item
		某次测量时，选用的双缝的间距为 $0.300\Umm$，测得屏与双缝间的距离为 $1.20\Um$，第 1 条暗条纹到第 4 条暗条纹之间的距离为 $7.56\Umm$。则所测单色光的波长为 \tkanswer[1.8cm]{630} $\Unm$（结果保留 3 位有效数字）。
	\end{enumerate}
	\onepicture[w=9cm, align=right]{figs/34e.png}
\end{enumerate}
%% answer:
\jdanswer{
\begin{enumerate}
	\item
	A

	\item
	B，$\frac{\Delta x\cdot d}{(n-1)l}$，630

\end{enumerate}
}
%% memo:
\memoanswer{
（1）由于细绳偏离竖直方向的角度很小，可将该模型视为单摆，单摆做简谐运动的周期公式为 $T=2\pi\sqrt{\frac{l}{g}}$，所以碰钉子后，摆长变为 $\frac{l}{4}$，则周期变为原来的 $\frac{1}{2}$，由机械能守恒定律可知小球一定能升到左侧同等高度处，由于被钉子挡住，由几何关系可知钉子的水平位移大小达不到初始水平位移的大小，由于从正向最大位移处计时，所以图像为余弦函数，故 A 正确。

（2）（i）干涉条纹间距 $\Delta x_0=\frac{l}{d}\lambda$，其中 $d$ 是双缝间距，$l$ 是屏到双缝的距离。要想增加从目镜中观察到的条纹数，就要减小 $\Delta x_0$，可以在 $d$ 不变的情况下，减小屏到双缝的距离 $l$；或在 $l$ 不变的情况下，使用间距更大的双缝，故 B 正确。

（ii）根据题意，相邻两暗纹间的距离为 $\Delta x_0=\frac{\Delta x}{n-1}$，再利用条纹间距公式 $\Delta x_0=\frac{l}{d}\lambda$ 得 $\lambda=\frac{\Delta x\cdot d}{(n-1)l}$。

（iii）把数据代入上式可得 $\lambda=\frac{7.56\times10^{-3}\times0.300\times10^{-3}}{3\times1.20}\Um=0.63\times10^{-6}\Um=630\Unm$。
}

\end{enumerate}

\end{document}
